\documentclass[11pt,a4paper]{report}
\usepackage{fontspec}
\setmainfont{Latin Modern Roman}
\setsansfont{Latin Modern Sans}
% Characters missing in Latin Modern Mono (the emoji and Chinese characters in
% examples/01_hello.s) are taken from macOS fonts, if they exist.
\directlua{luaotfload.add_fallback("monofallback",
  {"Arial Unicode MS:mode=harf;", "Apple Color Emoji:mode=harf;"})}
\setmonofont{Latin Modern Mono}[Scale=0.88,RawFeature={fallback=monofallback}]
\usepackage[margin=27mm]{geometry}
\usepackage[skip=\medskipamount,indent=0pt]{parskip}
\usepackage{amsmath,amssymb}
\usepackage{microtype}
\usepackage{xcolor}
\usepackage{graphicx}
\usepackage{listings}
\usepackage{fancyvrb}
\usepackage{etoolbox}
\usepackage{booktabs}
\usepackage{array}
\usepackage{longtable}
\usepackage{enumitem}
\usepackage[font=small,labelfont=bf]{caption}
\usepackage{tikz}
\usepackage[most]{tcolorbox}
\usetikzlibrary{arrows.meta,positioning,fit,calc,decorations.pathreplacing,patterns,
  circuits.logic.US,matrix}
\usepackage{xurl}
\usepackage{titlesec}
\titleformat{\chapter}[display]{\normalfont\bfseries}{\large\chaptertitlename\ \thechapter}{4pt}{\LARGE}
\titlespacing*{\chapter}{0pt}{0pt}{16pt}
\usepackage{hyperref}
\hypersetup{
  colorlinks=true,
  linkcolor=blue!50!black,
  urlcolor=blue!50!black,
  linktoc=all,
  bookmarksnumbered=true,
  bookmarksopen=true,
  bookmarksopenlevel=0,
  pdftitle={ULM on ICE: Building a Computer on an FPGA, Step by Step},
  pdfauthor={Michael C. Lehn}
}

% Start a section on the next page if it would get only a line or two here.
\usepackage{needspace}
\usepackage{placeins}
\let\plainsection\section
\renewcommand{\section}{\needspace{6.5\baselineskip}\plainsection}

\definecolor{codebg}{RGB}{247,248,250}
\definecolor{codekw}{RGB}{0,70,140}
\definecolor{codecomment}{RGB}{95,105,115}
\definecolor{termbg}{RGB}{240,242,238}
\lstdefinestyle{code}{
  basicstyle=\ttfamily\small,
  keywordstyle=\color{codekw},
  commentstyle=\color{codecomment},
  backgroundcolor=\color{codebg},
  frame=single,
  rulecolor=\color{black!15},
  columns=fullflexible,
  keepspaces=true,
  showstringspaces=false,
  breaklines=true,
  tabsize=8,
  aboveskip=6pt,
  belowskip=6pt
}
\lstdefinelanguage{SystemVerilog}[]{Verilog}{
  morekeywords={logic,always_comb,always_ff,always_latch,typedef,enum,struct,
    packed,interface,endinterface,modport,package,endpackage,import,unique,
    priority,int,bit,byte,shortint,longint,localparam,genvar,generate,
    endgenerate},
  sensitive=true
}
% Commands typed at a prompt and what they print.
\lstdefinestyle{term}{
  style=code,
  language={},
  basicstyle=\ttfamily\footnotesize,
  backgroundcolor=\color{termbg},
  breaklines=true,
  breakindent=2em,
  postbreak={}
}
% \endgraf drops the indentation of an enclosing list, so we restore it by hand.
\makeatletter
\lstnewenvironment{terminal}{\endgraf\noindent\hskip\@totalleftmargin\minipage{\linewidth}\lstset{style=term}}{\endminipage\endgraf}
\makeatother
\lstset{style=code,language=SystemVerilog}

% Files of the repository, given relative to its top directory.  Whole files
% may break across pages; \svpart keeps a part of a file on one page.
\newcommand{\svfile}[2][]{\lstinputlisting[language=SystemVerilog,title={\texttt{\detokenize{#2}}},#1]{../#2}}
\newcommand{\svpart}[2]{\par\noindent\begin{minipage}{\linewidth}\lstinputlisting[language=SystemVerilog,title={\texttt{\detokenize{#2}}},#1]{../#2}\end{minipage}\par}
\newcommand{\cppfile}[2][]{\lstinputlisting[language=C++,title={\texttt{\detokenize{#2}}},#1]{../#2}}
\newcommand{\cpppart}[2]{\par\noindent\begin{minipage}{\linewidth}\lstinputlisting[language=C++,title={\texttt{\detokenize{#2}}},#1]{../#2}\end{minipage}\par}
\newcommand{\makefilepart}[2]{\par\noindent\begin{minipage}{\linewidth}\lstinputlisting[language=make,title={\texttt{\detokenize{#2}}},#1]{../#2}\end{minipage}\par}
% A file shown verbatim, for files with characters outside of Latin-1, which
% listings would put in the wrong order (examples/01_hello.s).
\newcommand{\verbfile}[1]{\par\noindent\begin{tcolorbox}[colback=codebg,colframe=black!15,
  boxrule=0.4pt,sharp corners,left=3pt,right=3pt,top=2pt,bottom=2pt,
  title={\texttt{\detokenize{#1}}},coltitle=black,colbacktitle=white,
  fonttitle=\small,toptitle=0pt,bottomtitle=2pt,halign title=center]
  \VerbatimInput[fontsize=\small,tabsize=8]{../#1}\end{tcolorbox}\par}
% Output of the tools, captured in the directory output.
\newcommand{\outputfile}[1]{\par\noindent\begin{minipage}{\linewidth}\lstinputlisting[style=term,basicstyle=\ttfamily\scriptsize]{output/#1}\end{minipage}\par}

\setlist[description]{leftmargin=1.5em,itemsep=2pt,topsep=2pt,font=\normalfont\bfseries}
\setlist[itemize]{itemsep=2pt,topsep=2pt}
\setlist[enumerate]{itemsep=4pt,topsep=2pt}
\setlength{\emergencystretch}{1.5em}

% Something for the reader to do.
\definecolor{tryitbar}{RGB}{40,90,190}
\newcounter{exercise}[chapter]
\renewcommand{\theexercise}{\thechapter.\arabic{exercise}}
\newenvironment{exercise}
  {\par\medskip\refstepcounter{exercise}\noindent
   \begin{list}{}{\leftmargin=9pt\rightmargin=0pt\topsep=0pt}\item\relax
   \textcolor{tryitbar}{\textbf{Exercise~\theexercise.}}\ }
  {\end{list}\medskip}

% Solutions at the end of a chapter, in the order of the exercises.
\newcounter{solution}[chapter]
\newcommand{\solutions}{\clearpage\section{Solutions}\setcounter{solution}{0}}
\newcommand{\solution}{\par\medskip\refstepcounter{solution}\noindent
  \textcolor{tryitbar}{\textbf{Exercise~\thechapter.\arabic{solution}.}}\ }

% A pitfall.
\definecolor{warnbar}{RGB}{170,70,20}
\newenvironment{pitfall}
  {\par\medskip\noindent
   \begin{list}{}{\leftmargin=9pt\rightmargin=0pt\topsep=0pt}\item\relax
   \textcolor{warnbar}{\textbf{Pitfall.}}\ }
  {\end{list}\medskip}

% Exclusive or: a dotted \lor.
\newcommand{\xor}{\mathbin{\dot{\lor}}}

% Terms defined in the text.
\newcommand{\term}[1]{\emph{#1}}
% Inline code on a light gray background, as on GitHub.
\definecolor{inlinebg}{RGB}{236,238,241}
\tcbset{inlinecode/.style={on line,colback=inlinebg,colframe=inlinebg,frame empty,boxrule=0pt,
  arc=2pt,boxsep=0pt,left=2pt,right=2pt,top=1pt,bottom=1pt,before upper={\vphantom{(}}}}
% \smash: the background may reach into the line spacing, but never changes it.
\DeclareRobustCommand{\cc}[1]{\smash{\tcbox[inlinecode]{\ttfamily\detokenize{#1}}}}

% Figures.
\tikzset{
  box/.style={draw,rounded corners=2pt,align=center,font=\small},
  toolbox/.style={box,fill=blue!7,minimum height=9mm},
  filebox/.style={box,fill=black!5,minimum height=9mm},
  chipbox/.style={box,fill=green!10},
  arrow/.style={->,>=Stealth,thick,draw=black!65},
  biarrow/.style={<->,>=Stealth,thick,draw=black!65},
  note/.style={font=\footnotesize,align=center,text=black!70},
  wire/.style={thick,draw=black!80},
  % timing diagrams: \draw[signal] ... with high at y=0.5, low at y=0
  signal/.style={thick,draw=blue!60!black},
  clocksig/.style={thick,draw=black!70},
}
